\paragraph{分支指标像空间的坐标逆极限与全深度 prime-register 的无限秩障碍}
在定理 \ref{thm:xi-branch-register-strict-inverse-limit} 已把自然扩张的分支寄存器模型识别为有限深度截断的严格逆极限之后，还可把这一事实重写为像空间坐标本身的逆极限完备性；进一步，当每一层都存在真实分支时，忠实的全深度 prime-register 账本便不再是有限层 prime-slice 最优化的简单延长，而会在结构层上强制出现可数无穷个彼此独立的素数方向。

\begin{theorem}[分支指标像空间本身就是有限截断坐标系的规范逆极限]\label{thm:xi-branch-register-image-coordinate-inverse-limit}
\leanverified{paper\_xi\_branch\_register\_image\_coordinate\_inverse\_limit}
沿用定理 \ref{thm:killo-natural-extension-branch-register} 的记号。对每个 \(n\ge 1\)，定义前缀投影
\[
\mathrm{pr}_{\le n}:\widehat X\to X\times \NN^n,
\qquad
\mathrm{pr}_{\le n}\bigl(x,(a_1,a_2,\dots)\bigr):=(x,a_1,\dots,a_n),
\]
并记
\[
Y_n:=\mathrm{pr}_{\le n}(\widehat X)
=
\Bigl\{(x_0,a_1,\dots,a_n)\in X\times\NN^n:\ \exists\,\widehat x\in\widehat X,
\mathrm{pr}_{\le n}(\widehat x)=(x_0,a_1,\dots,a_n)\Bigr\}.
\]
定义截断映射
\[
\rho_{n+1,n}:Y_{n+1}\to Y_n,
\qquad
\rho_{n+1,n}(x_0,a_1,\dots,a_{n+1})=(x_0,a_1,\dots,a_n).
\]
则规范映射
\[
\Theta:\widehat X\to \varprojlim(Y_n,\rho_{n+1,n}),
\qquad
\Theta(\widehat x):=\bigl(\mathrm{pr}_{\le n}(\widehat x)\bigr)_{n\ge 1},
\]
为双射。

进一步，对每个 \(n\ge 1\) 定义有限深度右移插入系统
\[
\widehat T_n:Y_n\to Y_n,
\qquad
\widehat T_n(x_0,a_1,\dots,a_n):=
\begin{cases}
(T(x_0),\beta(x_0)),& n=1,\\
(T(x_0),\beta(x_0),a_1,\dots,a_{n-1}),& n\ge 2.
\end{cases}
\]
则在上述识别下，
\[
\widehat T=\varprojlim_n \widehat T_n.
\]
换言之，分支指标模型 \(\widehat X\) 本身就是有限深度右移插入系统的逆极限对象。
\end{theorem}
\begin{proof}
由定理 \ref{thm:xi-branch-register-finite-past-minimal-complete-prefix}，\(n\)-步过去的分支寄存器前缀
\[
(x_0,a_1,\dots,a_n)
\]
与唯一有限前史
\[
(x_0,x_{-1},\dots,x_{-n})\in X^{\leftarrow,n}
\]
一一对应。反过来，由 \(T\) 为满射，任意有限前史都可继续向后逐层选取原像，因而总能延拓为某个无限前史。故 \(Y_n\) 与定理 \ref{thm:xi-branch-register-strict-inverse-limit} 中的
\[
\widehat X_n^{\mathrm{br}}=\mathcal C_n(X^{\leftarrow,n})
\]
是同一对象，只是记号改写而已。

\(\Theta\) 的单射性是直接的：若两点在全部有限前缀上相同，则其 \(\widehat X\subset X\times\NN^\NN\) 中的全部坐标都相同，因而两点一致。

再证满射。任取一致族
\[
y^{(n)}=(x_0^{(n)},a_1^{(n)},\dots,a_n^{(n)})\in Y_n,
\qquad
\rho_{n+1,n}(y^{(n+1)})=y^{(n)}.
\]
由一致性可知 \(x_0^{(n)}\) 稳定，且对每个固定 \(j\ge 1\)，当 \(n\ge j\) 时 \(a_j^{(n)}\) 亦稳定。记稳定值为
\[
x_0,\qquad a_j\in\NN\quad (j\ge 1).
\]
由于每个 \(y^{(n)}\in Y_n\)，由定理 \ref{thm:xi-branch-register-finite-past-minimal-complete-prefix} 可恢复唯一有限前史
\[
(x_0,x_{-1}^{(n)},\dots,x_{-n}^{(n)})\in X^{\leftarrow,n}.
\]
而 \(\rho_{n+1,n}(y^{(n+1)})=y^{(n)}\) 又迫使这些有限前史在重叠部分完全一致，故它们定义出唯一无限前史
\[
(x_0,x_{-1},x_{-2},\dots)\in X^{\leftarrow}.
\]
将其送入定理 \ref{thm:killo-natural-extension-branch-register} 的编码映射，得到
\[
\widehat x:=\mathcal C(x_0,x_{-1},x_{-2},\dots)\in \widehat X,
\]
并满足
\[
\mathrm{pr}_{\le n}(\widehat x)=y^{(n)}\qquad(\forall n\ge 1).
\]
故 \(\Theta\) 为满射。

最后，由定理 \ref{thm:killo-natural-extension-branch-register} 的右移插入公式
\[
\widehat T(x,(a_1,a_2,\dots))=(T(x),(\beta(x),a_1,a_2,\dots)),
\]
对每个 \(n\) 取前缀便得到
\[
\mathrm{pr}_{\le n}\!\bigl(\widehat T(\widehat x)\bigr)
=
\widehat T_n\!\bigl(\mathrm{pr}_{\le n}(\widehat x)\bigr).
\]
因此 \(\Theta\) 将 \(\widehat T\) 精确识别为有限深度右移插入系统的逆极限动力学。证毕。
\end{proof}

\begin{theorem}[忠实的全深度 prime-register 实现必然具有无限素数秩，因而排斥任何有限生成加法账本]\label{thm:xi-infinite-depth-prime-ledger-infinite-rank-obstruction}
\leanverified{paper\_xi\_infinite\_depth\_prime\_ledger\_infinite\_rank\_obstruction}
设 \(T:X\twoheadrightarrow X\) 为有限纤维满射，并假设其逆向分支深度无界：对每个 \(N\ge 1\)，都存在长度 \(N\) 的前史片段
\[
x_{-N}\xrightarrow{\,T\,}x_{-N+1}\xrightarrow{\,T\,}\cdots\xrightarrow{\,T\,}x_{-1}\xrightarrow{\,T\,}x_0
\]
满足
\[
\#T^{-1}(x_{-j+1})\ge 2
\qquad(1\le j\le N).
\]
称一个全深度 prime-register 账本实现为忠实，若它对每个有限深度 \(N\) 的前缀截断都与定理 \ref{thm:xi-layered-prime-register-sharp-natural-extension} 的精确分层外置语义相容。

则任意忠实的全深度实现都必须使用无穷多个彼此独立的 prime slices。更具体地，若从每一层的有效 prime slice 中选取一枚代表素数
\[
q_j\in \mathcal P_j\qquad (j\ge 0),
\]
并定义自由交换幺半群
\[
M_\infty
:=
\left\{
\prod_{j\ge 0}q_j^{a_j}:\ a_j\in\ZZ_{\ge 0},\ a_j=0\ \text{仅有限多个例外}
\right\},
\]
则 \(M_\infty\) 在唯一分解意义下同构于 \(\NN_{\times}\)。

特别地，不存在任何有限生成交换群 \(L\) 使忠实的全深度 prime-register 账本能够严格同态化到
\[
(\ZZ\oplus L,+)
\]
这样的有限秩加法账本上。
\end{theorem}
\begin{proof}
第一部分正是定理 \ref{thm:xi-natural-extension-infinite-prime-register-necessity} 的内容：在逆向分支深度无界时，任何精确实现整个自然扩张的分层 prime-register 外置都不可能只使用有限多个 prime slices。结合上一定理，\(\widehat X\) 已被识别为全部有限深度前缀系统的逆极限，因此忠实的全深度账本若与每个有限前缀的精确恢复相容，就必须在每一层保留一个独立的有效 prime direction。

由 \((q_j)_{j\ge 0}\) 两两不同且都是素数，唯一分解给出幺半群同构
\[
\vartheta:\NN_{\times}\xrightarrow{\ \sim\ } M_\infty,
\]
其在素数生成元上由
\[
\vartheta(p_{j+1})=q_j
\qquad(j\ge 0)
\]
确定。

反设存在有限生成交换群 \(L\)，使忠实的全深度账本可严格同态化到 \((\ZZ\oplus L,+)\) 上。则其 restriction 在上述独立素数方向上给出一个幺半群单射
\[
M_\infty\hookrightarrow (\ZZ\oplus L,+).
\]
与 \(\vartheta\) 复合便得到
\[
\NN_{\times}\hookrightarrow (\ZZ\oplus L,+)
\]
的单射幺半群同态。
但 \((\ZZ\oplus L,+)\) 是有限生成交换可消去幺半群，这与定理 \ref{thm:killo-prime-freedom-non-finitizability} 矛盾；等价地，这也与定理 \ref{thm:killo-godel-compression-not-finite-rank-homologizable} 所排除的有限秩同态化账本机制相冲突。故此类有限生成加法账本不存在。证毕。
\end{proof}

\noindent
配合定理 \ref{thm:xi-terminal-fixed-freezing-tv-collapse}、定理 \ref{thm:xi-fixed-freezing-escort-bounded-observable-collapse} 以及推论 \ref{cor:xi-time-part57bb-freezing-groundstate-subset-equidistribution} 可见：自然扩张首先在分支指标像空间层面构成严格逆极限对象；冻结 escort 态随后在统计层面指数塌缩到最大纤维上的微正则分布；而一旦要求把这种全深度可逆账本忠实地外置为 prime-register，其结构层便不可避免地转入可数无穷素数秩，从而排斥一切有限秩加法线性化。

\endinput
